% Template created by Karol Kozioł (www.karol-koziol.net) for ShareLaTeX

\documentclass[a4paper,9pt]{extarticle}
\usepackage[utf8]{inputenc}
\usepackage[T1]{fontenc}
\usepackage{graphicx}
\usepackage{xcolor}
\usepackage{tikz}

\usepackage{amsmath,amssymb,textcomp}
\everymath{\displaystyle}

\usepackage{times}
\renewcommand\familydefault{\sfdefault}
\usepackage{tgheros}
\usepackage[defaultmono,scale=0.85]{droidmono}

\usepackage{multicol}
\setlength{\columnseprule}{0pt}
\setlength{\columnsep}{20.0pt}


\usepackage{geometry}
\geometry{
a4paper,
total={210mm,297mm},
left=10mm,right=10mm,top=10mm,bottom=15mm}

\linespread{1.3}


% custom title
\makeatletter
\renewcommand*{\maketitle}{%
\noindent
\begin{minipage}{0.65\textwidth}
\begin{tikzpicture}
\node[rectangle,rounded corners=6pt,inner sep=10pt,fill=blue!50!black,text width= 0.95\textwidth] {\color{white}\Huge \@title};
\end{tikzpicture}
\end{minipage}
\hfill

\bigskip\bigskip
}%
\makeatother

% custom section
\usepackage[explicit]{titlesec}
\newcommand*\sectionlabel{}
\titleformat{\section}
  {\gdef\sectionlabel{}
   \normalfont\sffamily\Large\bfseries\scshape}
  {\gdef\sectionlabel{\thesection\ }}{0pt}
  {
\noindent
\begin{tikzpicture}
\node[rectangle,rounded corners=3pt,inner sep=4pt,fill=blue!50!black,text width= 0.95\columnwidth] {\color{white}\sectionlabel#1};
\end{tikzpicture}
  }
\titlespacing*{\section}{0pt}{15pt}{10pt}


% custom footer
\usepackage{fancyhdr}
\makeatletter
\pagestyle{fancy}
\fancyhead{}
\fancyfoot[C]{\footnotesize \textcopyright\ \@date\ \ \@author}
\renewcommand{\headrulewidth}{0pt}
\renewcommand{\footrulewidth}{0pt}
\makeatother


\title{Cortex for bacterial genomics: quickstart}
\date{2015}



\begin{document}

\maketitle

\begin{multicols*}{2}


\section{Requirements}

Download Cortex from github (https://github.com/iqbal-lab/cortex)  - there are download instructions there.

You must have installed VCFtools (and have the entire directory, not just the binary executable), R and Stampy. R must be in your path. Add the following directories to PERL5LIB and PATH
\begin{verbatim}
export PERL5LIB= /path/cortex/scripts/analyse_variants/
    bioinf-perl/lib;/path/cortex/scripts/calling:
    /path/VCFTools_dir/perl:$PERL5LIB
export PATH = /path/cortex/scripts/analyse_variants/
   needleman_wunsch-0.3.0
\end{verbatim}
Also, make an INDEX file, mapping sample-identifiers to sequence data. 

\section{Independent workflow: motivation}
Do variant discovery independently for each sample (against a reference), then combine the callsets to make a single set of candidate sites (SNPs, indels, SVs), and then genotype all samples. 

We already have a script for this - run\_calls.pl - but at the genotyping step this combines all samples into one graph. This won't scale to thousands of samples, so we have a new pipeline. This cheatsheet shows how to run it on a single machine with many CPUs/cores, using GNU parallels.

\section{Independent workflow: how to run it}
This script will compile Cortex for you, make reference genome binaries, Stampy indexes - all the things that used to have to be done manually. It will also choose appropriate memory-use parameters
\begin{verbatim}
 perl cortex/scripts/calling/run_indep_wkflow_with_gnu_par.pl  
 --index INDEX  --ref_fa ref_genome.fa  
 --dir_for_ref_objects ref/ 
 --vcftools_dir ~/installed_apps/vcftools_0.1.9/ 
 --outdir results/ 
 --stampy_bin ~/installed_apps/stampy-1.0.23/stampy.py 
 --kmer 31 --procs 20 --prefix salmonella
\end{verbatim}
This starts with FASTQ (or BAM) files, and finishes with a single VCF with all samples genotyped at the same sites. 

\section{Troubleshooting}



\section{Segregating variants within our dataset (joint workflow)}
First build sample graphs as Step1 in the previous example. Then UNDETERMINED



\section{Pan-genome analysis}

To detect presence of a set of predefined genes (genes.fasta) among your samples
\begin{verbatim}
\end{verbatim}
To look at pan-genome graph of all samples and see which samples have which contigs, allowing you to stratify them by frequency or look for differentiating/segregating contigs.


\section{The End}

For further information:


\end{multicols*}

\end{document}
